\appendixsection{Möbius function}

Let \(n\) be an integer \(\geqslant 1\) . If \(n\) is divisible by the square of a prime number, we write \(\mu(n) = 0\) . If \(n\) is not divisible by the square of a prime number, we write \(\mu(n) = (-1)^k\) , where \(k\) is the number of prime divisors of \(n\) . The function \(\mu: \mathbf{N}^* \to \{-1, 0, 1\}\) thus defined is called the Möbius function.

Recall that given two integers \(n_1 \geqslant 1, n_2 \geqslant 1\) , we write \(n_1|n_2\) if \(n_1\) divides \(n_2\) .

PROPOSITION. (i) The function \(\mu\) is the unique mapping of \(\mathbf{N}^*\) into \(\mathbf{Z}\) such that \(\mu(1) = 1\) and

\[
\sum_ {d \mid n} \mu (d) = 0\tag{1}
\]

for every integer \(n > 1\) .

\begin{enumerate}
\def\labelenumi{(\roman{enumi})}
\setcounter{enumi}{1}
\tightlist
\item
  Let \(s\) and \(t\) be two mappings of \(\mathbf{N}^*\) into a commutative group written additively. In order that
\end{enumerate}

\[
s (n) = \sum_ {d | n} t (d) \quad \text { for   every   integer } n \geqslant 1,\tag{2}
\]

it is necessary and sufficient that

\[
t (n) = \sum_ {d \mid n} \mu (d) s \left(\frac {n}{d}\right) \quad for every integer n \geqslant 1.\tag{3}
\]

The uniqueness assertion in (i) is obvious, for (1) allows us to determine \(\mu(n)\) by induction on \(n\) . We show that the function \(\mu\) satisfies (1). Let \(n\) be an integer \(>1\) . Let P be the set of prime divisors of \(n\) and let \(n = \prod_{p \in \mathbb{P}} p^{\nu_p(n)}\) be the decomposition of \(n\) into prime factors. If \(d\) is a divisor of \(n\) , then \(\mu(d) = 0\) unless \(d\) is of the form \(\prod_{p \in \mathbb{H}} p\) , where H is a subset of P. Then

\[
\begin{array}{r l} \sum_ {d | n} \mu (d) & = \sum_ {\mathrm{H} \subset \mathrm{P}} (- 1) ^ {\mathrm{CardH}} \\ & = \sum_ {k = 0} ^ {\mathrm{CardP}} \left(\frac {\mathrm{CardP}}{k}\right) (- 1) ^ {k} = (1 - 1) ^ {\mathrm{CardP}} = 0. \end{array}
\]


Let \(s\) and \(t\) be two mapping of \(\mathbf{N}^*\) into a commutative group written additively. Let \(n \in \mathbf{N}^*\) . If (2) holds, then

\[
\begin{array}{r l} \sum_ {d | n} \mu (d) s \left(\frac {n}{d}\right) & = \sum_ {d | n} \mu (d) \sum_ {\delta \mid (n / d)} t (\delta) = \sum_ {d \delta | n} \mu (d) t (\delta) \\ & = \sum_ {\delta | n} t (\delta) \sum_ {d \mid (n / \delta)} \mu (d) = t (n). \end{array}
\]

Conversely, if (3) holds, then

\[
\sum_ {d \mid n} t (d) = \sum_ {d \mid n} \sum_ {\delta \mid d} \mu (\delta) s \left(\frac {d}{\delta}\right) = \sum_ {d \mid n} s (d) \sum_ {\delta \mid (n / d)} \mu (\delta) = s (n),
\]

which completes the proof.

Formula (3) is called the Möbius inversion formula.


